\documentclass[11pt]{article}

\usepackage[margin=1in]{geometry}
\usepackage{amsmath,amssymb,amsthm,mathtools,booktabs,tabularx}
\usepackage{microtype,xurl}
\usepackage{xcolor}
\usepackage[colorlinks=true,linkcolor=blue!55!black,
            citecolor=blue!55!black,urlcolor=blue!55!black]{hyperref}

\newtheorem{theorem}{Theorem}[section]
\newtheorem{lemma}[theorem]{Lemma}
\newtheorem{proposition}[theorem]{Proposition}
\theoremstyle{remark}
\newtheorem{remark}[theorem]{Remark}
\newcommand{\T}{\mathbb T}
\newcommand{\norm}[1]{\left\lVert#1\right\rVert}
\newcommand{\cert}[1]{\nolinkurl{#1}}
\newcolumntype{Y}{>{\raggedright\arraybackslash}X}

\title{A Compensated Green-Tail Counterfamily\\
\large for the Bateman--Turok Euclidean Torus}
\author{GPT-5.6.sol\thanks{OpenAI model and principal programme author.}
\and Asger Alstrup Palm\thanks{Programme orchestrator and corresponding human contact; Honorary Professor, DTU Compute, Technical University of Denmark. \texttt{asger@area9.dk}.}}
\date{Working draft, 17 August 2026}

\begin{document}
\maketitle

\begin{abstract}
The positive Euclidean lattice proposed by Bateman and Turok invites a
deterministic route to reconstruction: prove that the complete action
gradient cannot become softer than the free bilaplacian scale.  We formulate
that route as a scaled Polyak--\L ojasiewicz inequality on the isotropic
four-torus and disprove it.  An explicit positive, nonseparable,
polynomial-contrast family has action bounded above and away from zero while
its normalized gradient quotient satisfies
\(Q_n/\omega_{L_n}^{2}=O(n^{-2})\to0\).  The construction uses an exactly
mean-compensated critical source, the periodic graph Green operator, and a
small four-way tensor perturbation.  It realizes, rather than evades, the
previously proved necessities of residual escape, unweighted-curvature
flatness, height-cut current cancellation, and nonvanishing action.  The
result refutes one all-field deterministic proof architecture only.  It does
not decide the Witten quotient, Gibbs typicality, the interacting
\(H^{-1}\) moment, continuum reconstruction, Born/Krein interpretation, or
any Lorentzian claim.
\end{abstract}

\section{Question and exact boundary}

Bateman and Turok propose to control a higher-derivative Euclidean theory by
a positive lattice weight and a hidden ghost-parity structure
\cite{BatemanTurok2026}.  Before reflection positivity, continuum limits, or
operator reconstruction can be addressed, one may ask a more elementary
finite-volume question: is there a deterministic action-gradient floor at
the free spectral scale?

Let \(\T_L^4=(\mathbb Z/L\mathbb Z)^4\), with nearest-neighbor graph
Laplacian \(\Delta\), and let \(u:\T_L^4\to(0,\infty)\).  Define
\[
 r_x=\frac{\Delta u_x}{u_x},\qquad
 A(u)=\frac12\sum_x r_x^2,
 \qquad g=\nabla_{\log u}A,
 \qquad Q(u)=\frac{\norm{g}_2^2}{\norm{r}_2^2}.
\]
The free gap is
\[
 \omega_L=4\sin^2\!\left(\frac{\pi}{L}\right).
\]
The proposed deterministic route would require a volume-independent
\(c>0\) such that \(Q(u)\geq c\omega_L^2\) for every nonconstant positive
field.  Our result is the negation of precisely this statement.

\begin{theorem}[Compensated Green-tail counterfamily]
\label{thm:main}
There are explicit integers \(L_n\to\infty\) and positive nonseparable
fields \(u_n\) on \(\T_{L_n}^4\) for which constants \(0<c<C<\infty\),
independent of \(n\), satisfy
\[
 c\leq\norm{r(u_n)}_2^2\leq C,
 \qquad \frac{\max u_n}{\min u_n}=O(L_n^{5/12}),
 \qquad \frac{Q(u_n)}{\omega_{L_n}^2}=O(n^{-2})\longrightarrow0.
\]
Consequently no positive volume-independent constant gives the all-field
scaled inequality above.
\end{theorem}

The dependency tags of Theorem~\ref{thm:main} are exactly
\texttt{LOCAL-ALGEBRAIC} and \texttt{EUCLIDEAN-SPECTRAL}.  It is neither a
\texttt{REDUCED-MODE} nor a \texttt{LORENTZIAN-CAUSAL} theorem.

\section{The weighted-current identity}

The useful variable is not \(r\) by itself but
\[
 h_x=\frac{r_x}{u_x^2},\qquad c_{xy}=u_xu_y.
\]
Direct differentiation, retaining both orientations of every edge, gives
the exact identity
\begin{equation}
 g_x=\sum_{y\sim x}c_{xy}(h_y-h_x)=:(L_ch)_x.
 \label{eq:weighted-current}
\end{equation}
Thus the nonlinear action gradient is a weighted graph Laplacian applied to
an unweighted curvature.  This identity explains both the concentration
theorems and the counterfamily: \(r=u^2h\) can retain finite energy at large
field height while \(h\) becomes flat and the weighted currents cancel.

For the height set \(S_K=\{x:u_x>K\}\), internal currents cancel and
\[
 \Gamma_K:=\sum_{x\in S_K}g_x
\]
is exactly the net current across its boundary.  Pairing with the centered
cut indicator and using \(\omega_L\geq16/L^2\) gives
\begin{equation}
 \frac{Q}{\omega_L^2}\geq
 \frac{\Gamma_K^2}{4\pi^4\norm{r}_2^2}.
 \label{eq:cut}
\end{equation}
The counterfamily must therefore cancel every fixed-height net current; it
cannot be a single uncancelled spike.

\section{The construction}

Take the integer scales
\[
 L_n=n^{24},\qquad \lambda_n=n^{16},\qquad R_n=n^{21},
 \qquad \epsilon_n=\frac{\lambda_n}{R_n^2}=n^{-26}.
\]
Let \(d_L(x,0)\) denote periodic graph distance and set
\begin{equation}
 f_n(x)=\frac{8\lambda_n^3}
   {(\lambda_n^2+d_{L_n}(x,0)^2)^3}.
 \label{eq:source}
\end{equation}
There is a unique mean-zero solution
\begin{equation}
 -\Delta v_n=f_n-\overline f_n.
 \label{eq:green}
\end{equation}
Define
\[
 u_n^0=v_n-\min v_n+\epsilon_n.
\]
The subtraction of \(\overline f_n\) is load-bearing.  It distributes the
compensating negative source through the periodic Green tail; a local cutoff
would create an order-one shell and restore a large gradient.

To exclude any hidden product ansatz, put
\[
 p_n(x)=\prod_{a=1}^4\sin\!\left(\frac{2\pi x_a}{L_n}\right),
 \qquad u_n=u_n^0(1+n^{-1}p_n).
\]
The factor is positive.  Its four-way Fourier coefficient has nonzero pairing
with \(\log u_n\) and is orthogonal to every sum of one-coordinate functions;
hence \(u_n\) is genuinely nonseparable.

\section{Estimates}

The proof uses exact finite-difference inequalities and four uniform analytic
estimates.  The detailed rational witnesses and the independent finite
fixture are stored in the certificate named in
\S\ref{sec:reproduction}.

\begin{lemma}[Action does not vanish]
There are dimension-only \(c,C>0\) such that
\[
 c\leq\left\|\frac{\Delta u_n^0}{u_n^0}\right\|_2^2\leq C.
\]
\end{lemma}

The inner annuli of the critical source supply the lower bound.  Periodic
Green estimates, split dyadically at \(\lambda_n\) and \(R_n\), control the
tail and give the upper bound.  The lift \(\epsilon_n\) prevents a zero of
the denominator without erasing the inner contribution.

\begin{lemma}[Gradient estimate]
The base field satisfies
\[
 \norm{g(u_n^0)}_2^2\leq C\left(
 \lambda_n^{-8}+\frac{\lambda_n^4}{R_n^8}
 +\frac{R_n^4}{L_n^8}
 +\frac{\lambda_n^4}{R_n^4L_n^4}\right).
\]
\end{lemma}

The four terms respectively record the smoothed core, finite-radius tail,
periodic compensation, and core--tail interaction.  They are kept separate
in the certificate so that no asymptotic cancellation is assumed.

\begin{lemma}[Stable nonseparable perturbation]
\[
 \norm{r(u_n)-r(u_n^0)}_2=O(n^{-1}),\qquad
 \norm{g(u_n)-g(u_n^0)}_2=O(n^{-1}L_n^{-2}).
\]
\end{lemma}

Since \(\omega_L^2\geq256/L^4\), substitution of the four scales yields
\[
 \frac{Q(u_n)}{\omega_{L_n}^2}
 \leq C(n^{-32}+n^{-8}+n^{-12}+n^{-20}+n^{-2}).
\]
This proves Theorem~\ref{thm:main}.  Notice that the slowest term is the
deliberate nonseparability perturbation; the radial Green-tail base collapses
faster.

\section{How the proof search narrowed the geometry}

The final construction was not guessed in isolation.  A sequence of exact
results removed simple architectures and derived necessary properties of
any counterfamily.

\begin{center}
\begin{tabularx}{\textwidth}{@{}p{0.18\textwidth}YY@{}}
\toprule
Stage & What was established & What remained \\
\midrule
RF--83--85 & Cycle hierarchies collapse only relative to a diameter method;
one-phase torus pullbacks and independent tensor phases move away from the
free scale. & A genuinely coupled four-dimensional corrector. \\
RF--86--92 & Sparse-flow, top-band, dyadic and virial bounds progressively
exclude super-polynomial contrast and action density above \(272/29\). & A
bounded-action, polynomial-contrast sequence. \\
RF--93--94 & Collapse at positive action forces residual energy above every
fixed height, flat \(h=r/u^2\), and cancellation across every height cut. &
Compatibility of all three conditions. \\
RF--95 & A uniform small-action floor excludes \(A\to0\). & The
positive-action concentration branch. \\
RF--96 & The compensated Green tail realizes that branch and refutes the
all-field floor. & Probabilistic, interacting, continuum and reconstruction
questions. \\
\bottomrule
\end{tabularx}
\end{center}

This history matters logically.  RF--93--95 are not failed attempts at the
final theorem: they are necessary-condition theorems that the counterfamily
must and does satisfy.  Conversely, RF--96 does not invalidate the earlier
large-action or large-contrast lower bounds; it lives in their surviving
region.

\section{What the counterfamily changes}

The theorem removes a tempting shortcut.  Uniform deterministic coercivity
of the complete finite-volume score cannot by itself carry the free
bilaplacian gap through all positive field configurations.  Any successful
reconstruction argument must use information absent from that statement.
Three possibilities remain conceptually distinct:

\begin{enumerate}
\item a probabilistic theorem showing that Green-tail neighborhoods have
negligible Gibbs weight;
\item an interacting \(H^{-1}\) or Witten-form estimate whose negative or
nonlocal terms see more than \(\norm{\nabla A}^2/\norm{r}^2\);
\item a reconstruction mechanism that needs only typical or observable
sectors rather than an all-field deterministic floor.
\end{enumerate}

The first route requires an entropy-versus-action comparison for a positive
radius tube around the family, not just the weight of one configuration.  The
second requires the actual interacting operator, not an analogy with the
free gap.  The third must still establish reflection positivity or its stated
replacement, a probability rule, and a controlled continuum limit.

\section{Claim boundary}

The theorem does \emph{not} establish any of the following:
\begin{itemize}
\item failure of the full Bateman--Turok proposal or of its Witten quotient;
\item non-negligible Gibbs probability of the counterfamily;
\item failure of an interacting \(H^{-1}\) moment or spectral gap;
\item a continuum counterexample or obstruction to Osterwalder--Schrader
reconstruction;
\item a Born rule, Krein probability interpretation, positive physical state,
particle space, scattering or unitarity;
\item any claim about Lorentzian propagation, BV Green operators, Hadamard
states, renormalized products or a Lorentzian QME.
\end{itemize}

These exclusions are not cautionary prose appended after the fact.  They are
machine-readable negative claim flags in the certificate and claim map.

\section{Certification and reproduction}
\label{sec:reproduction}

The authoritative result is
\begin{center}
\path{reverse_physics/certificates/REVERSE_PHYSICS_BT_EUCLIDEAN_TORUS_GREEN_TAIL_COUNTERFAMILY_V1.json}.
\end{center}
It records the scale ledger, exact inequalities, dependency tags, controls,
input hashes, assumptions and exclusions.  Its verifier is independent of
the exploratory floating-point scout.  An exact rational \(4^4\)-torus
fixture reduces the source-centered geometry to 15 symmetry orbits while
retaining all 256 vertices; it checks the Green solve, positivity, residual,
weighted-current identity, nonseparability witness and exponent ledger.

The affected deterministic test chain is run by
\begin{verbatim}
python3 -m unittest discover -s reverse_physics/tests \
  -p 'test_bt_euclidean_torus_*.py'
python3 paper/generate_22_bateman_turok_torus_claim_map.py --check
python3 paper/verify_22_bateman_turok_torus_claim_map.py
\end{verbatim}
The paper claim map hash-pins RF--83 through RF--96.  Re-running a producer is
reproduction; the independent verifier is the separate audit rail.

\begin{thebibliography}{9}

\bibitem{BatemanTurok2026}
S.~Bateman and N.~Turok,
``Escape from Ostrogradsky via hidden ghost parity,'' 2026.
\url{https://arxiv.org/abs/2607.00096}.

\bibitem{ReverseFoundations2026}
OpenAI Codex and A.~A.~Palm,
``Reverse foundations of physics: separating physical postulates,
mathematical carriers, and logical strength,'' Paper~21 of this programme,
2026; \path{paper/21-reverse-foundations-of-physics.pdf}.

\end{thebibliography}
\end{document}
